\section{A precise route toward quasi-polynomial recognition}
\label{sec:research}

The new component predicate improves the value of each supplied normal
candidate.  A candidate may carry irrelevant closed or boundary-parallel
components, yet its useful disk is still detected.  The scheduler improves
one operation repeatedly used by that predicate.  To turn these gains into
a global complexity theorem, one must control how candidates and geometric
states are obtained, how many are visited, and how their bit sizes evolve.

\subsection{A conditional assembly statement}

The following formulation makes these obligations explicit.  It is a
research specification, not an assertion that its hypotheses have been
implemented.

\begin{proposition}[Conditional recognition through compressed candidates]
\label{prop:conditional-recognition}
Suppose an algorithm, on an $n$-crossing knot diagram $D$, constructs in
polynomial time a compact triangulation $E(D)$ and a checked correspondence
with its knot exterior.  Suppose a complete candidate procedure has the
following properties:
\begin{enumerate}
\item It visits at most $S(n)$ states, including unsuccessful branches, with
an encoded state-size bound $B(n)$.
\item Each transition, each normal-vector construction, and all bookkeeping
take time polynomial in $B(n)$, including all integer bit costs.
\item Its terminal acceptance condition is the presence of a compressing
disk among the components of a supplied admissible normal vector; a terminal
rejection certificate or an exhaustive termination theorem establishes
completeness of rejection.
\end{enumerate}
Then adding the component query gives an exact recognizer with time
$S(n)\operatorname{poly}(B(n))$, apart from the polynomial exterior
construction.  In particular, if $B(n)$ is polynomial and
$S(n)=2^{O((\log n)^a)}$ for a fixed $a\geq1$, the resulting bound is
quasi-polynomial.
\end{proposition}

\begin{proof}
The component query is polynomial in each encoded input.  Its positive
predicate is exactly the existence of a compressing-disk component, and
the verified exterior correspondence makes boundary compressibility
equivalent to the input being the unknot.  The assumed rejection theorem
rules out unvisited positive instances after termination.  Summing the
polynomial transition and query costs over all visited states gives the
bound.  No component multiplicity is expanded in this sum.
\end{proof}

The proposition also shows why a fast verifier alone is insufficient.
Normal coordinates of bounded bit size still range over a vast set; bounding
their encoding is not a bound on the number of candidates examined.  A
successful search heuristic does not supply the rejection theorem in the
third hypothesis.

Lackenby's hierarchy iteration expression $L(g+1)^L$ provides a more
geometric target \cite{Lackenby2026}.  If $g\leq n^c$ and all operations
and state encodings satisfy polynomial bounds, then
\[
 \log\bigl(L(g+1)^L\bigr)=O(\log L+L\log n).
\]
Thus $L=O(\log n)$ would give $2^{O((\log n)^2)}$, while
$L=O((\log n)^2)$ gives $2^{O((\log n)^3)}$.  The latter agrees with the
exponent displayed in the 2021 Oxford handout \cite{Lackenby2021}.
The expression becomes a runtime bound only after the geometric operations,
branch count, and representation sizes have been bounded constructively.
It is not enough to count hierarchy levels while allowing an exponential
triangulation expansion inside a level.

\subsection{Research questions with concrete completion criteria}

\paragraph{1. Certify the passage from a diagram to a compact exterior.}
Can a diagram-to-triangulation constructor emit a short, independently
checkable correspondence including the peripheral marking?  The immediate
deliverable should be a source-bound certificate whose verification accepts
exactly the claimed exterior, with a proved polynomial tetrahedron and bit
bound.  This closes the first missing logical link before the new component
predicate can affect the ordinary recognizer.

\paragraph{2. Find candidates without enumerating all admissible vectors.}
Which hierarchy or handle moves produce a complete family of useful normal
vectors while visiting only quasi-polynomially many states?  The component
predicate permits extra disconnected summands, potentially relaxing the
candidate specification.  A useful theorem would show that a cheaply
constructed larger surface must contain the desired disk, without requiring
enumeration of every fundamental or vertex normal surface.  Any such theorem
must cover unsuccessful inputs as well as successful examples.

\paragraph{3. Recover a selected component in compressed coordinates.}
Given a positive five-coordinate profile, can one produce one actual disk
component and its $7t$ normal-coordinate vector in polynomial bit time?
One candidate approach first filters with five weights, then replays the
same orbit trace with one coordinate for each normal-piece type and suitable
representative information.  Classical weighted AHT makes the coordinate
payload plausible, but selecting a particular component inside a repeated
profile group and certifying its relation to the source remain explicit
tasks.  Returning all exponentially many components is not a suitable
completion criterion.

\paragraph{4. Preserve three-dimensional information through cutting.}
Can a verified disk or hierarchy surface be cut out while retaining a
polynomial representation of the residual manifold, parallelity regions,
peripheral curves, and attaching maps?  Scalar weights or set-incidence
profiles do not contain cyclic order or three-dimensional gluing data.
The compressed cutting and parallelity framework in
\cite[Theorem~14.4]{Lackenby2026} is a natural target for constructive
implementation and bit-size accounting.  A complete step should include
both an output representation and a checker for its geometric meaning.

\paragraph{5. Obtain a useful global potential.}
What potential controls hierarchy length, candidate branching, coordinate
growth, and the work of unsuccessful reductions simultaneously?  Existing
local simplifications can decrease a word length while increasing another
representation.  A successful potential should yield a recurrence for total
work, not merely demonstrate progress in one favorable parameter.  The
parameters $g$ and $L$ should be connected to quantities computed from the
input, with every representation conversion included.

\paragraph{6. Realize or exclude the cubic family in surface geometry.}
The five-cycle family proves an asymptotic improvement for legal interval
systems.  Does a comparable family arise from valid normal coordinates on
compact torus-boundary triangulations, preferably from verified knot
exteriors?  Alternatively, do geometric arc systems satisfy a structural
restriction that rules it out and yields a stronger scheduling theorem?
Either outcome would clarify which worst cases matter to the intended
recognizer.  The present native meridian measurements show smaller constant
factor gains and do not answer this question.

\paragraph{7. Reduce the event queue's memory without losing its theorem.}
Can one find all potentially successful periodic mergers using
$o(k^2)$ stored candidates while retaining a proved $O(k^2)$ or better
arithmetic work bound?  A naive per-row minimum can become invalid when a
candidate is deleted and can silently restore repeated full scans.  Useful
approaches must specify how failed pairs are remembered, how generation
changes are propagated, and how all rescans are amortized.  The full
trace-order guarantee is valuable, but a different order is acceptable only
with a renewed trace-length and correctness analysis.

\paragraph{8. Tune adaptive scheduling using measurable work.}
The bounded prelude already protects easy first-pair successes from eager
queue initialization.  What inexpensive indicators predict when building
the queue will pay for its memory and arithmetic costs?  The preliminary
experiment exposed a separate periodic-index rebuild cost, leading to the
lazy first-pair path included here.  Future policies should retain a
mathematical test allowance and be evaluated with paired timings on easy,
native, adversarial, and large-bit inputs.  Universal speedup should not be
assumed from a worst-case improvement.

\paragraph{9. Stream weighted replay and share geometric preparation.}
Can the producer emit checked operations to an online weight accumulator,
while also saving a replayable certificate, to avoid retaining redundant
state?  Can repeated candidates reuse a validated triangulation and the
boundary basis without trusting stale source data?  The desired output is a
clear ownership and source-binding contract, together with a peak-memory
bound.  These changes can improve practical searches without changing the
global mathematical exponent.

\paragraph{10. Extend component features only when they answer a needed
question.}
Which constant-dimensional additive predicates suffice for subsequent
hierarchy choices?  Euler characteristic and two mod-two evaluations suffice
for the torus compressing-disk predicate, but they do not recover slopes,
boundary-circle partitions, or embeddings.  Additional integer or oriented
weights may recover some of this information; nonadditive information may
require a richer compressed structure.  Each added field should have an
explicit interpretation, bit bound, and independent test.

\paragraph{11. Understand the higher-genus boundary obstruction.}
For genus $h>1$, $2h$ mod-two coordinates detect a nonzero boundary homology
class but miss essential separating curves.  Therefore replacing two cycles
by $2h$ cycles does not extend the disk-essentiality theorem.  Can a
compressed simple-curve algorithm decide essentiality for a selected disk
boundary while preserving source markings?  A completion criterion must
handle separating curves, multiple boundary components, and their
interaction with the chosen component, rather than only aggregate homology.

\paragraph{12. Shrink and verify the trusted geometric core.}
The current trace verifier is independent of scheduling, whereas geometric
preparation and weight construction are shared.  A second implementation or
a formal proof should target local edge orientation transport, the
one-cell/one-corner allocation, boundary homology bases with repeated
incidences, and full-range weight transfer.  The fixed Regina fixtures and
literal finite oracles are useful regression evidence, not substitutes for
these proofs.  Formalization should be accompanied by encoding and bit-cost
lemmas so that correctness does not conceal expansion.

\subsection{Recommended next integration sequence}

The most useful immediate sequence is to integrate the proved scheduler and
component query, add a certified exterior constructor, and recover a single
positive component in normal coordinates.  Those steps produce an exact
geometric acceptance path with inspectable evidence.  In parallel, a
hierarchy prototype should instrument state count, bit growth, parallelity
compression, and unsuccessful branches.  Its purpose is to test and refine
the hypotheses of \cref{prop:conditional-recognition}; small experimental
state counts should not be reported as a quasi-polynomial proof.

The long-term mathematical bottleneck is the controlled global geometric
search.  The present work removes avoidable repeated merger tests and
enriches the compressed certificate available at each candidate.  It makes
that bottleneck more explicit and gives subsequent research a stronger,
tested local interface.
